\scanpage{8}
\mspage{- 55 -}
n\textsuperscript{o}357

[N.B. -- Ce n'est pas par amour des méthodes éculées et pisseuses que le rédacteur, dans la construction de la mesure de Haar, est revenu aux mesures à support \uline{fini} au lieu des mesures à support \uline{compact} des rédactions précédentes ; c'est en vertu du rasoir d'Occam (``entitates non sunt multiplicandae''), et pour pouvoir laisser reposer en paix dans la Remarque de la fin du Chap. VII, \handwritten{\S}\ 2, n\textsuperscript{o}2, la théorie générale de l'opération $\mu\ast f$, ou $\handwritten{\gamma}(\mu)f$, lorsqu'on ne s'est pas donné de jauge ; ce n'est évidemment pas difficile (cf. loc. cit), mais ça paraît bien superflu \ldots]

On écrira $\handwritten{\unsure{\mathcal{K}|\mathscr{K}}}$, $\handwritten{\mathcal{K}}(K)$, etc. au lieu de $\handwritten{\mathcal{K}}(G)$, $\handwritten{\mathcal{K}}(G,K)$, etc. ; on notera $\handwritten{\mathcal{K}}_+$ l'ensemble des fonctions $f\mathrel{\handwritten{\in}}\handwritten{\mathcal{K}}$ à valeurs $\mathrel{\handwritten{\geq}}0$, $\handwritten{\mathcal{K}}^{\ast}_{+}$ l'ensemble des fonctions $g\mathrel{\handwritten{\in}}\handwritten{\mathcal{K}}_+$ autres que $0$, et $\handwritten{\mathcal{K}}^{\ast}_{+}(K)$, si $K$ est une partie compacte de $G$, l'ensemble des fonctions $g\mathrel{\handwritten{\in}}\handwritten{\mathcal{K}}^{\ast}_{+}$ à support contenu dans $K$. Pour $f\mathrel{\handwritten{\in}}\handwritten{\mathcal{K}}_+$, $g\mathrel{\handwritten{\in}}\handwritten{\mathcal{K}}^{\ast}_{+}$, désignons par $A(f,g)$ l'ensemble des mesures positives $\mu\mathrel{\handwritten{\ne}}0$, à support fini, telles que $f\mathrel{\handwritten{\leq}}\handwritten{\gamma}(\mu)g$. \uline{Cet ensemble n'est pas vide}. En effet, soient $K$ le support de $f$, et $U$ un ouvert non vide tel que $m=\inf_U(g)>0$ ; soit $M=\sup_K(f)$. On peut recouvrir $K$ par un nombre fini d'ensembles $s_i U$ ; alors la mesure $(M/m)\handwritten{\sum}\varepsilon_{s_i}$ est dans $A(f,g)$. Posons maintenant :
\[
(f:g)=\inf_{\mu\mathrel{\handwritten{\in}} A(f,g)}\lVert\mu\rVert\ .
\]

Il est clair qu'on a $0\mathrel{\handwritten{\leq}}\quad(f:g)<+\handwritten{\infty}$. De plus, ce symbole a les propriétés suivantes :

(i) \quad $(\handwritten{\gamma}(\varepsilon_s)f:g)=(f:g)$ quel que soit $s\mathrel{\handwritten{\in}}G$.

(ii) \quad $(cf:g)=c.(f:g)$ quel que soit $c\mathrel{\handwritten{\geq}}0$.

(iii) \quad $f\mathrel{\handwritten{\leq}} f'$ entraîne $(f:g)\mathrel{\handwritten{\leq}}(f':g)$.

(iv) \quad $(f:g)\mathrel{\handwritten{\geq}}(\sup f)/(\sup g)$.

(v) \quad $(f:h)\mathrel{\handwritten{\leq}}(f:g).(g:h)$ pour $g\mathrel{\handwritten{\in}}\handwritten{\mathcal{K}}^{\ast}_{+}$, $h\mathrel{\handwritten{\in}}\handwritten{\mathcal{K}}^{\ast}_{+}$.

(vi) \quad si $f$, $f_{\unsure{o|0}}$, $g$ sont dans $\handwritten{\mathcal{K}}^{\ast}_{+}$, on a :
% archive sheet number: 7

